\documentclass[10pt,a4paper]{amsart}
\usepackage[margin=19mm]{geometry}
\usepackage{amsmath,amssymb,mathtools}
\usepackage[scr=rsfs,cal=euler]{mathalfa}
\usepackage{xfrac}
\usepackage[unicode,hidelinks,bookmarks=false]{hyperref}
\newcommand{\ie}{i.e.\ }
\newcommand{\cf}{cf.\ }
\newcommand{\resp}{respectively\ }
\newcommand{\Subsection}{\subsection}
\providecommand{\nobreakdash}{\nobreak\hbox{-}}
\setlength{\emergencystretch}{2em}
\multlinegap0pt
\usepackage{amssymb}
\usepackage{tikz}
\usepackage{xcolor}
\newtheorem{corollary}{Corollary}
\newtheorem{theorem}{Theorem}
\newcommand{\ssym}[0]{spectrally symmetric }
\newcommand{\ssymend}[0]{spectrally symmetric}
\title{Spectrally symmetric orientations of graphs}
\author{Saieed Akbari, Jonathan Aloni, Maxwell Levit, Bojan Mohar, Steven Xia}
\date{Source: arXiv:2512.08049v2 (CC BY 4.0)}
\begin{document}
\maketitle

\noindent\emph{Source and licence.} Spectrally symmetric orientations of graphs. Version arXiv:2512.08049v2 (CC BY 4.0), distributed by arXiv under the Creative Commons Attribution 4.0 International licence, http://creativecommons.org/licenses/by/4.0/, as recorded in the arXiv licence metadata for this exact version and shown on the abstract page https://arxiv.org/abs/2512.08049v2. Full licence notice: \texttt{licenses/arxiv-2512.08049-CC-BY-4.0.txt}.\par\smallskip

\noindent\textit{Editorial scope.} Counted unit: the article's theorem Thm:10n (source0.tex lines 484--485). The article's complete original proof of it is delivered with it (source0.tex lines 487--504, one \texttt{proof} environment of 18 lines). No mathematics is written, repaired or re-derived: the statement and the proof are the article's own lines, in the article's own notation, and no retained line is substituted or re-flowed.  Retained uncounted as the source-local context the proof needs: lines 477--480.  Nothing was omitted from any retained span, so the package carries no omission record.  No retained line contains a citation, so the wrapper carries no bibliography and no bibliography entry.  No retained line contains a figure environment, an image-inclusion command or drawing code, so no image is delivered and the package declares no asset dependency.  Editorial formatting only, in addition: an AMS \texttt{amsart} preamble that restates the article's own declarations and macros used by the retained lines, namely the environments \texttt{corollary}, \texttt{theorem} on separate counters, together with the standard packages those lines need.  The retained environments print in this document as Corollary 1, Theorem 1 in order of appearance, whereas the article prints the same items with the numbering of its own full document.  The complete native source of the article as distributed in the arXiv e-print for this exact version is bundled unchanged as \texttt{sources/190625/source0.tex}.
\begin{corollary}\label{Cor:t<6e}
If $L$ is a \ssym line graph then
\[t(L)\le6|E(L)|.\]  
\end{corollary}

\begin{theorem}\label{Thm:10n}
    If the line graph $L=L(G)$ is \ssymend, then $|E(G)|\le10|V(G)|$.\end{theorem}

\begin{proof}

Let $G$ be a graph with $n$ vertices. Suppose $L$ is \ssymend. Then using Corollary \ref{Cor:t<6e} and counting triangles and edges of $L$ from the degrees of $G$ we have

\[t(G)+\sum_{i=1}^n\binom{d_i}{3}\le6\sum_{i=1}^n\binom{d_i}{2}.\]
After simplification and ignoring the $t(G)$ term we have
\begin{align}\sum\limits_{i=1}^nd_i^3-21d_i^2+20d_i<0\hspace{10pt} \label{eq_cubic}
\end{align}
So it suffices to show that  (\ref{eq_cubic}) fails when $|E(G)|> 10n$. Let $\alpha$ denote the average degree of $G$. Suppose $|E(G)|> 10n$ and thus $\alpha>20$. Then

\begin{align}\sum\limits_{i=1}^n(d_i^3-21d_i^2+20d_i)=&\sum_{i=1}^n((d_i-7)^3-127d_i+343)\nonumber\\ 
\geq&\sum_{i=1}^n((\alpha-7)^3-127\alpha+343)\label{eq_alpha_cubic}\\
=&\sum\limits_{i=1}^n(\alpha^3-21\alpha^2+20\alpha)\nonumber\\
=&~n(\alpha-20)(\alpha-1)\alpha \nonumber\\
>&~0.\nonumber\end{align}

The inequality on Line (\ref{eq_alpha_cubic}) follows because, for $\alpha>14$, $\sum_{i=1}^n(d_i-7)^3$ is minimized when all $d_i$ are equal.
\end{proof}

\end{document}
